\documentclass[11pt]{article}
\usepackage{../shared/luxcover}
\usepackage{amsmath,amssymb}
\usepackage{hyperref}
\usepackage{booktabs}
\usepackage{enumitem}
\usepackage{xcolor}

\title{Lux Quantum Swap}
\author{Zach Kelling\\\textit{Lux Industries, Inc.}\\\texttt{research@lux.network}}
\date{May 2024}

\begin{document}

\luxcoverpage

\begin{abstract}
QuantumSwap is a quantum resistant, completely private, gas-efficient, multi-chain, hyper liquid AMM (Automated Market Maker) protocol designed to operate on the Lux Network. The purpose is to facilitate RWA (real world asset) token swaps using customizable bonding curves that support both fungible and non-fungible tokens. QuantumSwap is the trading engine that facilitates the trade of assets on the Lux Market and Lux Exchange.

QuantumSwap supports both BTC and EVM based blockchains for universal cross chain liquidity.
\end{abstract}

\tableofcontents
\newpage


QuantumSwap supports both BTC and EVM based blockchains for universal cross chain liquidity.


Liquidity providers (LPs) can deposit assets natively on Bitcoin and Ethereum (as well as other EVM chains like Avalanche and Polygon) into a Quantum Vault and enable quantum secure trading and storage of nearly any token. Quantum Safe Assets (QSAs) may be traded with zero gas fees natively on Lux, with complete privacy, and we will create the tools to do programmatic signing and trading across blockchains.


QuantumSwap allows trading between liquidity pools using both ERC-20 and ERC-721 tokens, unlike traditional order book systems. The team is exploring support for Cosmos, Sui, and other blockchains to make QuantumSwap a universal, quantum-secure liquidity protocol.


\subsection{Evolution of DeFi Liquidity}

\subsection{Benefits for Quantum Swap AMM LP's (Liquidity Providers):}

LPs can efficiently exchange real-world asset ownership in the form of ERC-721 tokens, ensuring high liquidity.


Homomorphic Encryption protects AMM order book transactions, preventing frontrunning.


LPs enjoy up to 5000x capital efficiency compared to Uniswap v2, earning higher returns while leveraging real-world assets.


Capital efficiency allows low-slippage trading, outperforming centralized exchanges and stablecoin-focused AMMs.


LPs can increase exposure to preferred assets and reduce downside risk while programmatically trading based on factors of their choosing.


By adding liquidity above or below market price, LPs can execute fee-earning limit orders along a smooth curve.


\subsection{RWA (Real World Asset) Tokens}

Real World Asset (RWA) Tokens are backed by real-world assets like gold, silver, or uranium. These NFTs maintain 100\% or more value backed by the underlying asset (and the asset's global markets) and may be redeemed using a physical asset redemption contract as any traditional commodity or asset. RWA Tokens are non-fungible, non-divisible, and free from ownership transfer restrictions. Lux Quantum Swap's standard in post quantum encryption ensures that all Lux-issued RWAs are natively post-quantum safe.


RWA Tokens can be used as collateral for users staking RWA tokens to create price-stable currencies that earn fees from platform trading activities.


\subsection{Real World Assets (RWA) as Collateral}

The functionality of a RWA token is employed when users stake them to create price stability, and then LPs are rewarded with a fee for the yield earned as other users on the platform trade these assets.


For someone to be able to consistently trade LUXG (Gold Backed NFT), there needs to be a set amount of gold on the exchange ready to be traded. Furthermore, the issuer (selling) LUXG must have a corresponding quantity of physical GOLD as collateral. To guard against the GOLD/USD kind of volatility, Quantum Swap requires that the issuer of LUXG provide 1:1 the reserve of gold necessary to cover every single LUXG issued.


\subsection{Quantum Safe Liquidity Pool}

A QSA liquidity pool is a smart contract that allows you to instantly swap between two different kinds of real world assets. On QuantumSwap, one type of liquidity pool is an LUXG<>LUXU pool, which means that anyone holding physical Gold can instantly swap those Gold backed NFT's (LUXG) for Uranium Backed NFT (LUXU) or vice versa.


Ideally, a pool contains some amount of both assets, enabling users to swap back and forth between them. However, it's also possible to create a pool with one asset, meaning that users will only be able to buy that asset from the pool as needed.


Importantly, liquidity pools can be comprised of fungible or non-fungible tokens or a mix of both, enabling significantly more liquidity for NFTs, which generally do not have on-chain liquidity, and making it cheap and gas efficient to make complicated trades between FT -> NFT -> FT and so on.


Each LP is secured by post quantum encryption ensuring that liquidity providers can safely deposit and collect fees from trading without worrying about possible quantum attack. Programmatic trading allows LPs to set parameters and manage pool governance.


\subsection{Step by Step Guide to Becoming a LP (Liquidity Provider):}

Step 1: Liquidity providers can create pools backed by single FTs or NFTs, pairs or FT <> NFT liquidity pools. They choose whether they would like to buy or sell NFTs (or both) and specify a starting price and bonding curve parameters.


Step 2: Users can then buy tokens or sell tokens into these pools. Every time an asset is bought or sold, the price to buy or sell another item changes for the pool based on its bonding curve.


Step 3: At any time, liquidity providers can change the parameters of their pool or withdraw their Real World Assets.


\subsection{Price Stable Currencies backed by Real World Assets}

One of the major impediments to holding any asset that is not backed by real world value is unpredictable volatility. The more volatile an asset's relative price is to other commonly-held assets (BTC VS GOLD), the more risk there is in holding it. If volatility is sufficiently high, such assets become too risky to be used in everyday commercial activities. Crypto-currencies, being a relatively new asset class, are significantly more volatile than their non-digital counterparts. BTC/USD pricing's volatility slowing over time, but still settles on or around 5 to 6\%. Compare that to a more conventional asset volatility, such as Gold/USD exchange rate, which is only around 0.5\%. With recent world activities, especially black swan events, volatility and price spreads have only increased for commodities like gold and silver.


One way to mitigate market volatility and risk is to have price-stable assets in the exchange that are collateralized against a relatively less-volatile asset, like GOLD. This asset acts as a stable anchor against which all other crypto-assets could be valued, and each unit of such an asset will always give a predictable return. Quantum Swap has built in price-stable assets. One such asset is a uranium-backed price stable currency that is designed to be equivalent to the TLV (total locked value) of the RWA Tokens staked into the DAO governing pool rewards and staking incentives. These reward and incentive models mint Price Stable currencies. The idea of such currencies is not new. BitShares, for example, has such currencies for trading in its ecosystem. However, pegged currencies within the BitShares ecosystem suffer from two distinct challenges. Pegged Currencies are too volatile (not backed by real world value) as they are collateralized only by BTS (Bitshares own token) which is another relatively volatile crypto-asset.


There is too little issuance of the pegged currencies as there are no incentives. In Quantum Swap's design, the collateral is based on real world assets such as Gold and gold, a far more stable foundation of value. Quantum Swap's built-in interest rates also give dynamically adjusted incentives for RWA tokens holders to issue price stable currency LUXG or LUXU.


\subsection{Incentives}

Having LUXG sounds like a great incentive to hold it, but why would anyone want to issue price stable currencies? That's where interest comes in. For every price stable currency of LUXG issued, the holder of the Price Stable Currency needs to pay the issuer of the physical gold asset an interest rate. Think of it as a fee paid to someone who's providing this big pool of collateral as a cushion against market volatilities. Since Quantum Swap is an AMM (automated market maker), it is guaranteed that One LUXG will always translate into exactly One Troy oz of Physical GOLD. This is due to the Liquidity Pool's bonding curve that automatically adjusts interest rates to make sure that issuers are properly incentivized to issue more or less of LUXG in order to balance out supply and demand. For example, if demand for LUXG is rising, that means 1 LUXG will trade for more than the price of an OZ of GOLD. When this happens, incentives become increased to the supply of LUXG in order to bring that ratio back to a 1 to 1 trading ratio by increasing the interest rate. When the reverse happens, it's desirable to lower the interest rate to reduce the supply.


\subsection{Bonding Curve for RWA Backed Price Stable Currencies}

Quantum Swap uses LUXG vs USD real time premium to generate the current interest rate. The annualized interest rate moves from 0\\% to a 15\\% cap, while the premium goes from 0\% up.


In Quantum Swap initial iteration, the interest rate is calculated in the following manner:


b\_LUX = \$ GOLD oz value traded for 1 LUXG on Quantum Swap


b\_market = \$ GOLD oz value being traded for 1 USD on a collection of major exchanges


premium = (b\_LUX - b\_market)/b\_market


Setting the interest rate:


if (premium = 0\%): interest\_rate = 0\%


if (0\% < premium ≤ 2\%): interest\_rate = 2\%


if (2\% < premium ≤ 15\%): interest\_rate = premium


if (premium > 15\%): interest\_rate = 15\%


The initial 2\% is set to create a threshold to help trigger action on the part of the issuer - users tend to understand a fixed threshold better and are more likely to trade with one. Quantum Swap also caps the interest rate at 15\% to make sure that the holder of the asset won't feel insecure about a potentially limitless interest rate that continuously rises.


The AMM bonding curve then automatically calculates the interest rate based on the average premium.


\subsection{Multi-Chain Teleport Protocol}

independently developed blockchains can natively send data (primarily tokens) directly to one another. The majority of cross-chain bridges are developed by independent third parties, and depending on the protocols they interface with, they range greatly in maturity and security, potentially compromising their security and resulting in inconsistent user experiences.


The Quantum Swap teleport protocol utilizes a trustless relayer to transmit data across a zero-knowledge powered zkBridge, ensuring full post-quantum secure Lamport and lattice-based encryption.


\subsection{Current Modern Encryption Techniques are Compromised}

All of the most prevalent solutions used to secure sensitive information are based on encryption. However, encryption is not the only line of defense as there are numerous attacks that exist. Some of these include attacks on encryption keys, cryptographic attacks, corruption or integrity attacks, insider attacks, stolen ciphertext attacks, ransomware attacks, and data destruction. Moreover, many widely-used and well-established cryptographic algorithms, such as RSA-based encryption/signatures and elliptic curve cryptography, are not quantum safe.


\subsection{Encryption presents several vulnerabilities that make these attacks possible:}
\begin{itemize}
  \item First, a set of encryption keys (typically containing just one key) is used to decrypt an encrypted message. Therefore, encryption addresses the problem of preserving data confidentiality by creating the new problem of securing the set of encryption keys.
  \item Second, the security of encryption is considered largely in the computational setting, meaning it is based on cryptographic hardness assumptions. If an attacker does not have access to the set of encryption keys, the underlying security of the encryption primitive depends on the hardness of a mathematical problem, such as the factoring of a product of two large prime numbers. As hard as the problem may be, computers can always find a solution. For example, though the difficulty of factoring the product of two prime numbers increases as a function of the number of decimal digits of that product, a computer can always find a solution .
  \item Third, traditional encryption can only encrypt data when it is stored or sent between Information-nodes but not when it is processed by a node. In such a case, the data needs to be decrypted first, processed in plaintext, and then encrypted again, opening an important attack vector.
\end{itemize}

\subsection{Privacy and Homomorphic Encryption}

Homomorphic Encryption (HE) allows data to be processed while in encrypted form, addressing the third vulnerability described above. However, HE generally incurs high computational costs when applied to general Turingcomplete computations, and it still exhibits the other vulnerabilities present in standard encryption.


\subsection{Shamir's Secret Sharing (SSS)}

The following is a summary of the advantages of Lux ZChain when compared


to traditional encryption:

\begin{itemize}
  \item Encrypted data is one hack away, whereas ZChain data is T +1 hacks away. For the former, an attacker is required only to compromise the security of a single system to access the encrypted data. In contrast, since the blinding factors masking the data under ITS in ZChain are not in any one particular place, an attacker is required to hack a large proportion of all the nodes (which are themselves secured by ITS) on the ZChain Bridge to be able to access its data. Therefore, attacks in SMPC (secure multi party computation) networks (and also in ZChain) such as Sybil attacks, collusion between nodes, and bribery, typically involve controlling a large (T + 1 or more) group of nodes. These attacks also apply to other decentralized networks, such as public blockchains that dont use a zero knowledge file system.
  \item Encryption relies on cryptographic assumptions and can be broken. ZChain relies

on ITS and is crypto analytically unbreakable.

  \item Typically, encrypted data cannot be processed. It needs to be decrypted first, which opens an important attack vector (HE is an exception to this, but as addressed above, it tends to be computationally expensive and relies on cryptographic assumptions).

In ZChain, data can be transferred, stored, and processed while being


kept perfectly confidential and private.

  \item In traditional encryption, key management is typically a complex task that involves securing the generation, distribution, and storage of secret keys. ZChain solves this by allowing nodes to generate, distribute, and store shares of the blinding factors (which are equivalent to a key), and these are protected by the ITS properties of LSS.

We thus believe that ZChain will render traditional encryption-based networks obsolete. Once a secret is stored in ZChain, users will be able to authenticate and retrieve data and authorize processing as needed. Multi-Factor Authentication (MFA) will be used, combining biometrics, multi-signature schemes, device information, geographic location, and potentially many other verification factors adopted by the network and selected by the individual user. Due to the nature of the previously described blind computation of transformed information, nodes in the ZChain Network will perform this key security and gating function without any knowledge of what they are computing.

\end{itemize}

\subsection{ZChain (Zero Knowledge Bridge)}

ZChain is a privacy preserving subnetwork of Lux Network that utilizes zero knowledge proofs to allow for private transactions. Zero knowledge proof file systems and allow for a person to prove they know something without revealing what that something is. In the context of ZChain, this means that a person can prove that they own a certain asset without revealing what that asset is.


Zero knowledge bridges bring privacy to DeFi by ensuring that data remains confidential between parties. zero knowledge circuits allow for the verifiable computation of encrypted data, while zkVM provides a secure execution environment for smart contracts.


ZChain is a decentralized network that enables the private exchange of data and value. By using ZChain, Quantum Swap can keep data confidential, preventing fraud and protecting users' privacy.


\subsection{Post Quantum Safe}

In cryptography, a Lamport signature or Lamport one-time signature scheme is a method for constructing a digital signature. Lamport signatures can be built from any cryptographically secure one-way function; usually a cryptographic hash function is used.


Although the potential development of quantum computers threatens the security of many common forms of cryptography such as RSA, it is believed that Lamport signatures with large hash functions would still be secure in that event. Each Lamport key can only be used to sign a single message. However, combined with a hash tree, a single key can be used for many messages, making this a fairly efficient digital signature scheme.


\subsection{Gas free meta transactions}

Quantum Swap's ERC-2771 approach is a standardized technique for sending ``Gasless'' transactions. Meta transaction aware recipient contracts only rely on a small trusted forwarder contract for their security. This contract verifies the signature and nonce of the original sender. Quantum Swap provides a default implementation of Forwarder which verifies the signature and forwards the call to the recipient smart contract.


Meta transactions are transactions that have been authorized by a Transaction Signer and relayed by an untrusted third party that pays for the gas (the Gas Relay).


\subsection{Hyper Liquid QSA AMM (Automated Market Maker)}

Quantum Swap provides the infrastructure required to create quantum safe token liquidity pools. LPs (Liquidity Providers) can use capital efficiency to increase their exposure to preferred assets, reduce their downside risk, and execute trades with low slippage.


Furthermore, LPs can approximate a fee-earning limit order by adding liquidity to a price range above or below the market price along a smooth bonding curve.


\subsection{Bonding Curve Algorithms For QSA Pricing}

QSA pools use a bonding curve to determine the relative price at which one asset is traded for another. The more an asset is bought from the pool (increased buy Pressure), the more expensive that asset becomes. Conversely, the more an asset is sold to the pool (increased sell pressure), the cheaper that asset becomes.


A bonding curve is an algorithmic formula that defines the relationship between an asset's trading price and its supply. Bonding curves are a key feature of automated market makers since they are used to algorithmically adjust the price of an asset as it is traded against another asset.


To determine pricing, each LUXQSAPair is associated with a specific bonding curve set by the LP. At present, there are three choices: LinearCurve, ExponentialCurve, and XYKCurve. In the future, customized bonding curve contracts may be created with Zap Protocol and whitelisted for use with LUXQSAPair.


Immediately after a user trades with a pair, the bonding curve is consulted to determine what the new price of the pair will be.


Technical note: If the bonding curves are intended to be pure, i.e. they do not modify the state of pairs that call them. The actual logic for input/output validation and price updating happens via the LUXRWAPair contract.


\subsubsection{Linear Bonding Curve}

The linear curve performs an additive operation to update the price. delta is assumed to be set properly by the LP to be the same precision of the pair's underlying token. The price of an NFT is increased by a flat amount (called delta) every time an item is bought from the pool.


If the pair has just sold an NFT by giving out an NFT and receiving tokens, the next price it will quote to sell NFTs at will be delta more. Conversely, if the pair has just bought an NFT by giving out tokens and receiving an NFT, the next price it will quote to purchase NFTs at will be delta less.


With a linear bonding curve, Conversely, the price of the NFT is decreased by that same flat amount every time an item is sold to the pool.


For example, a liquidity provider may create an NFT<> ETH pool with a Start Price of 1 ETH and a delta of 0.1 ETH. Assuming they provide enough liquidity, the price of an NFT will increase to 1.1 ETH after one item is purchased from the pool. After a second item is purchased, the price will increase to 1.2 ETH, and so on and so forth. At any point, if an NFT is sold to the pool, the price will decrease by 0.1 ETH.


\subsubsection{Exponential Bonding Curve}

With an exponential bonding curve, the price of an RWA is increased by a certain percentage (also called delta) every time an item is bought from the pool. Conversely, the price of the RWA is decreased equivalently every time an item is sold to the pool.


To calculate the equivalent decrease, convert the percentage to a decimal index (e.g. for 50\%, the index would be 1.5) and divide the price by this number.


For example, a liquidity provider may create an NFT<>ETH pool with a Start Price of 2 ETH and a delta of 50\%. Assuming they provide enough liquidity, the price of an NFT will increase to 2 + 50\% = 3 ETH after one item is purchased from the pool. After a second item is purchased, the price will increase to 3 + 50\% = 4.5 ETH, and so on and so forth. At any point, if an NFT is sold to the pool, the price will be divided by 1.5.


If the pair has just sold an NFT by giving out an NFT and receiving tokens, the next price it will quote to sell NFTs at will be multiplicatively delta more. Conversely, if the pair has just bought an NFT by giving out tokens and receiving an NFT, the next price it will quote to purchase NFTs at will be multiplicatively delta less.


\subsubsection{XYK Bonding Curve}

With an XYK curve, the price of an NFT is adjusted every time an item is bought from or sold to the pool, such that the product of two virtual reserves remains constant after every trade. These virtual reserves correspond to the number and value of NFTs the pool will buy or sell.


An additional concentration parameter allows liquidity providers to adjust (i.e. tighten or loosen) XYK curves.


The XYK curve adjusts prices such that the product (multiplication) of two virtual reserves remains constant after every trade. At pool creation,the virtual reserves are set as follows:


balance: the number of tokens being bought or sold (in the case of trade pools, whichever is greater) plus one


tokenBalance: the number of tokens being bought or sold (or whichever is greater) multiplied by the start price


To conform to the IPair interface, RWABalance is stored as delta and tokenBalance is stored as spotPrice.


The net price (exclusive of pool and protocol fees) for an XYK pair to sell ``X amount of'' NFTs is inputValueWithoutFee = (x * tokenBalance) / (Balance - x). Conversely, the net price for an XYK pair to buy X NFTs is outputValueWithoutFee = (x * tokenBalance) / (Balance + x).


The net price (exclusive of pool and protocol fees) for an XYK pair to sell ``X amount of'' NFTs is inputValueWithoutFee = (x * tokenBalance) / (RWABalance - x). Conversely, the net price for an XYK pair to buy X NFTs is outputValueWithoutFee = (x * tokenBalance) / (RWABalance + x).


Immediately after every trade, the virtual reserves are updated:


Pair sells: nftBalance = nftBalance - x and tokenBalance = tokenBalance + outputValueWithoutFee


Pair buys: nftBalance = nftBalance + x and tokenBalance = tokenBalance = tokenBalance - outputValueWithoutFee


\subsection{Determining Asset Spot Price}

In addition to modifying delta to change the pair's price reactivity, it is important to understand how the spotPrice variable in LUXRWAPair behaves.


The spot price refers to the instantaneous price of selling 1 NFT to the pair. The instantaneous price of buying 1 NFT from the pair is set to be the spotPrice adjusted upwards by 1 unit of delta with respect to the pair's bonding curve.


For example, say that we have a Trade LUXRWAPair for ETH with a spotPrice of 1 ETH, and a linear curve with a delta of 0.1 ETH. (Assume fee is 0.) Then a user selling 1 NFT to the pair would receive 1 ETH, whereas a user purchasing 1 NFT from the pair would have to send (1 + 0.1) = 1.1 ETH.


In other words, the price to purchase an NFT from a pair will always be delta greater (either additively or multiplicatively) than the price to sell an NFT to the pair.


Pair managers who are manually adjusting spotPrice should keep this in mind.


(Note that for simplicity, some examples in this documentation may use "spot price" to also refer to the instantaneous buy price.)


\subsection{Pricing For Multi-Swaps}

If a user buys or sells multiple NFT in one swap transaction, the spotPrice will update by delta for each NFT bought or sold.


For example, say that we have a Trade LUXRWAPair for ETH with a spotPrice of 1 ETH, and a linear curve with a delta of 0.1 ETH. (Assume fee is 0.)


If a user sells 5 NFT to this pair, they will receive:


1 ETH for the first NFT


0.9 ETH for the second NFT


0.8 ETH for the third NFT


0.7 ETH for the fourth NFT


0.6 ETH for the fifth NFT


At the end of the multi-swap, the new spotPrice will be set to 0.5 ETH.


\begin{table}[htbp]
\centering
\begin{tabular}{l|l|l}
\toprule
2013 & Bitshares & Decentralized Pure Order Books \\
\midrule
2016 & 0x Protocol & First gas efficient meta transactions \\
2018 & Uniswap V1 & First ever LP Tokens,  
First AMM To Allow ETH<>ERC20 \\
2020 & Uniswap V2 & Wrapped ETH, ERC20<>ERC20 Pairs, Oracle Price Feeds, Flash Swaps \\
2021 & Uniswap V3 & Concentrated liquidity, Range Orders, New fee tiers \\
2023 & QuantumSwap & First Quantum AMM. Gas-efficient, multi-chain, and hyper-liquid protocol that supports all fungible and non-fungible token standards. It features customizable bonding curves and comprehensive quantum security for trading RWA tokens. \\
\bottomrule
\end{tabular}
\end{table}

\end{document}
